Identify a symmetric matrix with its \(s=d(d+1)/2\) upper-triangular entries. The set \(\mathcal K_{a,b}\) is closed and bounded in \(\mathbb R^s\), hence compact. We give the diagonal construction of the limiting law. For every integer \(r\ge1\), make a finite Borel partition \(\mathcal P_r\) of \(\mathcal K_{a,b}\) into sets of Euclidean diameter at most \(2^{-r}\), with \(\mathcal P_{r+1}\) refining \(\mathcal P_r\). Such partitions are obtained recursively by intersecting each cell with a finite \(2^{-r-2}\)-ball cover and disjointifying the intersections in a fixed order.

There are only countably many cells in \(\bigcup_r\mathcal P_r\). Since every cell probability belongs to \([0,1]\), successive applications of Bolzano--Weierstrass followed by diagonal selection give a subsequence, still denoted \((n_k)\), for which
\[
 \ell(C)=\lim_{k\to\infty}P\{U_{n_k}\in C\}
 \tag{1}
\]
exists for every partition cell \(C\). The limits are nonnegative, sum to one on each level, and are consistent under refinement. They therefore define a probability distribution on infinite nested cell paths: at a cell of positive mass choose a child with conditional probability \(\ell(C')/\ell(C)\), and assign arbitrary child probabilities after a zero-mass cell. Choose a point \(u_C\in C\) in every nonempty cell. Along every path of positive probability, the representatives are Cauchy because a level-\(r+1\) representative and all later representatives lie in the same level-\(r\) cell and hence are at distance at most \(2^{-r}\). Compactness gives a limit in \(\mathcal K_{a,b}\). Let \(U\) be this limit and let \(\mu\) be its law. Consistency makes the path probability of a level-
\(r\) cell equal to \(\ell(C)\).

Let \(f\) be continuous. Compactness makes it uniformly continuous; write
\[
 \omega_f(h)=\sup\{|f(u)-f(v)|:u,v\in\mathcal K_{a,b},\ \|u-v\|_2\le h\},
 \qquad \omega_f(h)\downarrow0. \tag{2}
\]
For one representative \(u_C\) per cell,
\[
 \left|E[f(U_{n_k})]-\sum_{C\in\mathcal P_r}f(u_C)P\{U_{n_k}\in C\}\right|
 \le\omega_f(2^{-r}), \tag{3}
\]
and the same inequality holds with \(U\) and \(\ell(C)\). First let \(k\to\infty\) in the finite sum using (1), and then let \(r\to\infty\) using (2)--(3). This proves \(E[f(U_{n_k})]\to E[f(U)]=\int f\,d\mu\), as required.